% ============================================================
% Chapter 6: Experimental Design
%   Section: Construction of the Custom Validation Set
%     Subsection: Size and Stratification
% Included from chapters/06-experimental-design/03-custom-validation-set/custom-validation-set.tex
% ============================================================

\subsection{Size and Stratification}
\label{subsec:custom-size}

The set targets \textbf{150 claims}, with 100 as the minimum below
which per-stratum counts stop being informative. With 150 claims, a
95\% confidence interval on an accuracy near 50\% is approximately
$\pm 8$ percentage points ($1.96\sqrt{0.25/150}\approx0.080$), which
is sufficient to separate the models compared in
Chapter~\ref{ch:results} when they differ substantially, and
insufficient to rank models that differ by a few points. That limit
is stated here rather than discovered later.

The claims are stratified along two axes, so that neither a single
verdict nor a single subject dominates the metrics:

\begin{itemize}
    \item \textbf{Verdict}: the five labels of the system's own
    \texttt{Verdict} enumeration (Table~\ref{tab:custom-labels}).
    \item \textbf{Topic}: each claim carries one of the 23 topics of
    the system's topic classifier (\texttt{src/config/topics.py}).
    Because 23 topics over 150 claims leaves six or seven claims per
    topic, too few to balance on, balance is sought over the five
    \emph{topic groups} into which those topics are already
    organised: society, science, environment, culture and health.
\end{itemize}

Five verdicts by five groups gives 25 cells, with a target of six
claims each (30 per verdict and 30 per group). Language (English or
Spanish) and claim type (factual, numerical or interpretive) are
recorded and reported, but not targeted.

The targets are a goal, not a quota. Positive-news articles are
predominantly accurate, so the \texttt{FALSE} and \texttt{MISLEADING}
cells are expected to be the hardest to fill from the outlet's own
articles. A cell that remains short is reported with its real count
rather than padded with claims from outside the domain, and the
evaluation reports macro-averaged $F_1$ alongside accuracy for that
reason: under class imbalance, accuracy rewards a system that
predicts the majority verdict.
